\documentclass[../../../include/open-logic-section]{subfiles}

\begin{document}

\olfileid[pt]{sth}{ord-arithmetic}{using-addition}
\olsection{Usando Adição Ordinal}

Com a adição nos ordinais, podemos calcular explicitamente os postos
de vários conjuntos, no sentido de \olref[spine][rank]{defnsetrank}:

\begin{lem}\ollabel{rankcomputation}
Se $\setrank{A} = \alpha$ e $\setrank{B} = \beta$, então:
\begin{enumerate}
	\item\ollabel{exrankpow} $\setrank{\Pow{A}} = \alpha\ordplus 1$
	\item\ollabel{exrankpair} $\setrank{\{A, B\}} = \max(\alpha,
	\beta) \ordplus 1$
	\item\ollabel{exrankcup} $\setrank{A \cup B} = \max(\alpha,
	\beta)$
	\item\ollabel{exranktuple} $\setrank{\tuple{A,B}} = \max(\alpha,
	\beta) \ordplus  2$
	\item\ollabel{exranktimes} $\setrank{A \times B} \leq \max(\alpha,
	\beta) \ordplus  2$
	\item\ollabel{exrankunion} $\setrank{\bigcup A} = \alpha$ quando
	$\alpha$ é vazio ou limite; $\setrank{\bigcup A} = \gamma$ quando
	$\alpha = \gamma\ordplus 1$
\end{enumerate}
\end{lem}

\begin{proof}
Ao longo de todo o raciocínio, invocamos repetidamente \olref[spine][rank]{ranksupstrict}.

\emph{\olref{exrankpow}.} Se $x \subseteq A$ então $\setrank{x} \leq
\setrank{A}$. Logo $\setrank{\Pow{A}} \leq \alpha \ordplus  1$. Como em particular $A
\in \Pow{A}$, tem-se $\setrank{\Pow{A}} = \alpha \ordplus  1$.

\emph{\olref{exrankpair}.} Por \olref[spine][rank]{ranksupstrict}

\emph{\olref{exrankcup}.} Por \olref[spine][rank]{ranksupstrict}.

\emph{\olref{exranktuple}.} Por \olref{exrankpair}, duas vezes.

\emph{\olref{exranktimes}.} Note-se que $A \times B \subseteq
\Pow{\Pow{A \cup B}}$, e invoque \olref{exranktuple}.

\emph{\olref{exrankunion}.} Se $\alpha = \gamma\ordplus 1$, existe
algum $c \in A$ com $\setrank{c} = \gamma$, e nenhum !!{element} de $A$
tem posto superior; logo $\setrank{\bigcup A} = \gamma$. Se $\alpha$ é um
ordinal limite, então $A$ tem !!{element}s com posto arbitrariamente próximo
(mas estritamente inferior) de $\alpha$, pelo que $\bigcup A$ também tem
!!{element}s com posto arbitrariamente próximo (mas estritamente inferior) de
$\alpha$, e portanto $\setrank{\bigcup A} = \alpha$.
\end{proof}
\noindent
Deixamos como exercício mostrar por que \olref{exranktimes} envolve uma
\emph{des}igualdade.

\begin{prob}
Exiba conjuntos $A$ e $B$ tais que $\setrank{A \times B}=
\max(\setrank{A}, \setrank{B})$. Exiba conjuntos $A$ e $B$ tais que
$\setrank{A \times B}=\max(\setrank{A}, \setrank{B}) \ordplus  2$. São
possíveis outros postos?
\end{prob}

Estamos também agora em condições de mostrar que várias noções razoáveis
do que possa significar descrever um ordinal como ``finito'' ou ``infinito'' coincidem:

\begin{lem}\ollabel{ordinfinitycharacter}
Para qualquer ordinal $\alpha$, as seguintes afirmações são equivalentes:
\begin{enumerate}
	\item\ollabel{ord:notinomega} $\alpha\notin \omega$, isto é,
	$\alpha$ não é um número natural
	\item\ollabel{ord:omegaplus} $\omega \leq \alpha$
	\item\ollabel{ord:oneplus} $1 \ordplus  \alpha = \alpha$
	%\alpha \approx \alpha \ordplus 1$
	\item\ollabel{ord:plusone} $\alpha \approx \alpha\ordplus 1$,
	isto é, $\alpha$ e $\alpha\ordplus 1$ são equipotentes
	\item\ollabel{ord:infinite} $\alpha$ é infinito de Dedekind
	\end{enumerate}
\end{lem}
\noindent
Temos assim cinco formas provavelmente equivalentes de precisar o que exige que um ordinal seja (in)finito.

\begin{proof}
\emph{\olref{ord:notinomega} $\Rightarrow$ \olref{ord:omegaplus}.} Por
tricotomia.

\emph{\olref{ord:omegaplus} $\Rightarrow$ \olref{ord:oneplus}.} Fixe
$\alpha \geq \omega$. Por indução transfinita, existe um ordinal $\gamma$
mínimo (possivelmente $0$) tal que existe um ordinal limite $\beta$ com
$\alpha = \beta \ordplus \gamma$. Ora:
\[
	1 \ordplus \alpha =
	1 \ordplus (\beta \ordplus \gamma) =
	(1 \ordplus \beta) \ordplus \gamma =
	\supstrict_{\delta < \beta} (1 \ordplus  \delta) \ordplus  \gamma =
	\beta \ordplus  \gamma =
	\alpha.
\]
\emph{\olref{ord:oneplus} $\Rightarrow$ \olref{ord:plusone}.} Existe
claramente !!a{bijection} $f \colon (\alpha \disjointsum 1) \to (1
\disjointsum \alpha)$. Se $1 \ordplus \alpha = \alpha$, existe um
isomorfismo $g \colon (1 \disjointsum \alpha) \to \alpha$. Considere-se
$\comp{f}{g}$.

\emph{\olref{ord:plusone} $\Rightarrow$ \olref{ord:infinite}.} Se
$\alpha \approx \alpha \ordplus 1$, existe !!a{bijection} $f \colon
(\alpha \disjointsum 1) \to \alpha$. Defina $g(\gamma) = f(\gamma, 0)$
para cada $\gamma < \alpha$; esta !!{injection} mostra que $\alpha$
é infinito de Dedekind, pois $f(0,1) \in \alpha \setminus \ran{g}$.

\emph{\olref{ord:infinite} $\Rightarrow$ \olref{ord:notinomega}.} Isto
é \olref[z][infinity-again]{naturalnumbersarentinfinite}.
\end{proof}

\end{document}
